\documentclass[a4paper,10pt]{ltjsarticle}

\usepackage[top=23mm,bottom=24mm,left=21mm,right=21mm]{geometry}
\usepackage{luatexja-fontspec}
\usepackage{amsmath,amssymb,amsthm,mathtools}
\usepackage{booktabs,longtable,tabularx,array,multirow}
\usepackage[dvipsnames]{xcolor}
\usepackage{enumitem}
\usepackage{hyperref}

\IfFontExistsTF{Yu Mincho}{\setmainjfont{Yu Mincho}}{}
\IfFontExistsTF{Yu Gothic}{\setsansjfont{Yu Gothic}}{}

\hypersetup{
  colorlinks=true,
  linkcolor=MidnightBlue,
  citecolor=MidnightBlue,
  urlcolor=MidnightBlue,
  pdftitle={ABS/APS上のG2, Consistency, Loeb, fixed points, cut elimination zoo},
  pdfauthor={My-Reserch-Project},
  pdfsubject={Abstract bimodal/provability structures and residuated APS},
  pdfkeywords={APS, ABS, G2, FG2, Loeb, Rosser, Jeroslow, fixed point, residuation}
}

\setlength{\parindent}{1em}
\setlength{\parskip}{0.2em}
\setlist[itemize]{leftmargin=2.1em,itemsep=0.12em}
\setlist[enumerate]{leftmargin=2.1em,itemsep=0.12em}
\renewcommand{\arraystretch}{1.15}

\newtheorem{theorem}{定理}
\newtheorem{proposition}[theorem]{命題}
\newtheorem{lemma}[theorem]{補題}
\newtheorem{definition}[theorem]{定義}
\newtheorem{remark}[theorem]{注意}
\newtheorem{question}[theorem]{問題}

\newcommand{\ABS}{\mathrm{ABS}}
\newcommand{\APS}{\mathrm{APS}}
\newcommand{\RAPS}{\mathrm{RAPS}}
\newcommand{\Gtwo}{\mathrm{G2}}
\newcommand{\FGtwo}{\mathrm{FG2}}
\newcommand{\nFGtwo}{\mathrm{nFG2}}
\newcommand{\Con}{\mathrm{Con}}
\newcommand{\Ros}{\mathrm{Ros}}
\newcommand{\Boxop}{\Box}
\newcommand{\Btx}{\boxtimes}
\newcommand{\Dia}{\Diamond}
\newcommand{\PrT}{\mathrm{Pr}_T}
\newcommand{\Fml}{\mathrm{Fml}}
\newcommand{\Fix}{\mathrm{Fix}}
\newcommand{\True}{\mathrm{Tr}}
\newcommand{\cut}{\mathrm{cut}}
\newcommand{\cf}{\mathrm{cf}}
\newcommand{\id}{\mathrm{id}}
\newcommand{\resl}{\backslash}
\newcommand{\resr}{/}
\newcommand{\Top}{T}
\newcommand{\Bot}{\bot}
\newcommand{\Repo}{\texttt{My-Reserch-Project}}

\title{ABS/APS上のG2・Consistency・Löb・不動点・cut elimination zoo}
\author{My-Reserch-Project}
\date{2026年6月11日}

\begin{document}
\maketitle

\begin{abstract}
本稿は，二重様相代数構造 \(\ABS\)，抽象的証明可能性構造 \(\APS\)，および剰余付き
\(\APS\)（本稿では \(\RAPS\) と略す）上で，さまざまな無矛盾性 statement，
\(\Gtwo/\FGtwo\)，Löb と formalized Löb，cut elimination，および
Gödel・Henkin・Santa Claus・Jeroslow 型不動点がどのように絡むかを整理するための
研究ノートである。結論を一行で言えば，\(\Gtwo\) は order-theoretic な外部原理，
\(\FGtwo\) はその内部化，Löb は consistency orbit の付着性，Rosser はその脱付着性，
cut elimination は剰余付き構造での資源合成の安定性として現れる。有限反例モデルは
\texttt{code/models/examples} と \texttt{artifacts/reports} に機械検証済みの形で残し，
本稿ではその読解表を与える。
\end{abstract}

\tableofcontents

\section{記法と基本構造}

\begin{definition}[ABS, APS, RAPS]
本稿で \(\ABS\) と呼ぶものは，少なくとも
\[
S=(L,\le,\Top,\Bot,\Boxop,\Btx)
\]
からなる二重様相的な順序構造である。\(\Boxop\) は証明可能性，
\(\Btx\) は反証可能性または refutability を表す。順序は
\[
\Top\le x \quad\text{を「\(x\) は証明される」},\qquad
x\le\Bot \quad\text{を「\(x\) は反証される」}
\]
と読む。

\(\APS\) は \(\ABS\) に，典型的には次の A1--A4 を課したものである。
\[
\tag{A1}
x\le y\Rightarrow \Boxop x\le\Boxop y,\qquad
x\le y\Rightarrow \Btx y\le\Btx x,
\]
\[
\tag{A2}
\Top\le \Btx\Bot,
\]
\[
\tag{A3}
x\le\Boxop y,\ x\le\Btx y\Rightarrow x\le\Btx\Top,
\]
\[
\tag{A4}
\Btx x\le\Boxop\Btx x.
\]
ただしこのリポジトリの finite zoo では，A1--A4 の一部だけを満たす
preAPS も明示的に扱う。

\(\RAPS\) は，さらに資源合成と剰余
\[
(L,\le,\otimes,\mathbf 1,\resl,\resr)
\]
を持ち，
\[
a\otimes b\le c
\quad\Longleftrightarrow\quad
b\le a\resl c
\quad\Longleftrightarrow\quad
a\le c\resr b
\]
を満たす \(\APS\) または preAPS である。
\end{definition}

\begin{remark}[top/bottom discipline]
有限モデル検索では，\(\Top\) が最大元，\(\Bot\) が最小元であるとは限らない。
特に bottom discipline
\[
\forall x\in L,\quad \Bot\le x
\]
は追加条件として扱う。これは ex falso 型の弱化であり，剰余付きモデルを作るときに
決定的な役割を持つ。
\end{remark}

\begin{definition}[\(\Gtwo,\FGtwo,\nFGtwo\)]
整合性軌道の第一要素を
\[
d:=\Btx\Top
\]
と書く。このとき
\[
\Gtwo(S):\quad d\le\Bot\Rightarrow \Top\le\Bot,
\]
\[
\FGtwo(S):\quad \Btx d=\Btx\Btx\Top\le \Btx\Top=d,
\]
\[
\nFGtwo(k):\quad \Btx^{k+1}\Top\le \Btx^k\Top\qquad(k\ge1)
\]
と定義する。通常の \(\FGtwo\) は \(\nFGtwo(1)\) である。
\end{definition}

\begin{remark}
非崩壊モデル，すなわち \(\Top\not\le\Bot\) では，
\(\Gtwo\) は \(d\not\le\Bot\) と同値になりうる。この場合の \(\Gtwo\) は
「反例がまだ出ていない」という意味で vacuous であり，\(\FGtwo\) や不動点存在とは
分離して読む必要がある。
\end{remark}

\section{Consistency statement の階層}

以下では \(T\) を算術理論，\(\PrT(x)\) を証明可能性述語，
\[
\Boxop\varphi:=\PrT(\ulcorner\varphi\urcorner),\qquad
\Dia\varphi:=\neg\Boxop\neg\varphi
\]
とする。ABS/APS では，これらの statement は単一の命題というより，
どの構造まで内部化できるかを測るためのテスト族である。

\begin{definition}[主な consistency statement]
本稿では次を用いる。
\[
\Con_T^L:=\neg\PrT(\ulcorner\Bot\urcorner)
\]
を Löb-style consistency と呼ぶ。
\[
\Con_T^{EG}:=\exists x\in\Fml\ \neg\PrT(x)
\]
を external Gödel-style consistency と呼ぶ。これは数としてコードされた式全体を
外側から量化する形である。
\[
\Con_T^{IG}:=\exists \varphi\ \neg\PrT(\ulcorner\varphi\urcorner)
\]
を internal/formula-level Gödel-style consistency と呼ぶ。実際に APS 内で扱うには，
式コード・有限結合・可算結合のどれを許すかを別途指定する必要がある。
本稿では両者をまとめて \(\Con_T^G\) と書く場合がある。

Strong/Hilbert 型は
\[
\Con_T^S:=\forall\varphi\,
\bigl(\PrT(\ulcorner\varphi\urcorner)
\to
\neg\PrT(\ulcorner\neg\varphi\urcorner)\bigr),
\]
Hilbert--Bernays 型は
\[
\Con_T^H:=\forall x\in\Fml\,
\neg\bigl(\PrT(x)\land\PrT(\mathrm{neg}(x))\bigr)
\]
である。Rosser 型 \(\Ros\) は，Rosser proof predicate による無矛盾性，
または rule
\[
\neg\varphi\quad\Longrightarrow\quad \neg\Boxop\varphi
\]
として読む。
\end{definition}

\begin{proposition}[基本的な強さの目安]
標準的な証明可能性述語と通常の算術化のもとでは，おおまかに
\[
\Con_T^G
\ <\
\Con_T^L
\ <\
\Ros
\ <\
\Con_T^S
\ <\
\Con_T^H
\]
という順に強くなる。ここで \(<\) は「右側がより強い consistency statement」
という目安であり，具体的な導出には \(E,M,C,D3,CB\) などの条件が関与する。
\end{proposition}

\begin{center}
\begin{tabularx}{\linewidth}{>{\raggedright\arraybackslash}p{0.20\linewidth}
>{\raggedright\arraybackslash}p{0.34\linewidth}
>{\raggedright\arraybackslash}X}
\toprule
statement & 典型的な形 & 読み方 \\
\midrule
\(\Con_T^L\) & \(\neg\Boxop\Bot\) & 単一の矛盾定数が証明されない。Löb/G2 の標準入口。\\
\(\Con_T^{IG}\) & \(\exists\varphi\,\neg\Boxop\varphi\) & 証明されない式がある。内部化には式空間の扱いが必要。\\
\(\Con_T^{EG}\) & \(\exists x\in\Fml\,\neg\PrT(x)\) & コード化された式全体への外部量化。\\
\(\Con_T^S\) & \(\forall\varphi(\Boxop\varphi\to\Dia\varphi)\) & 証明された式は同時に反証されない。\\
\(\Con_T^H\) & \(\forall x\,\neg(\PrT(x)\land\PrT(\neg x))\) & Hilbert--Bernays 型の全式無矛盾性。\\
\(\Ros\) & Rosser consistency/rule & 反証側の「早い証明」を比較する。Löb から脱付着する余地を持つ。\\
\bottomrule
\end{tabularx}
\end{center}

\subsection{G2 の算術的 variant}

既存ノート \path{research/notes/g2-zoo-arithmetic.md} の比較表を，証明可能性条件の
観点で読み直すと次のようになる。

\begin{longtable}{>{\raggedright\arraybackslash}p{0.24\linewidth}
>{\raggedright\arraybackslash}p{0.32\linewidth}
>{\raggedright\arraybackslash}p{0.31\linewidth}}
\toprule
結論 & 十分条件 & コメント \\
\midrule
\endfirsthead
\toprule
結論 & 十分条件 & コメント \\
\midrule
\endhead
\(T\nvdash\Con_T^L\) &
\(\Sigma_1\)-predicate \(+\ D2+D3\) &
標準 HBL/Löb 経路。Gödel 文 \(\pi\leftrightarrow\neg\Boxop\pi\) を使う。\\
\(T\nvdash\Con_T^L\) &
\(\Sigma_1\)-predicate \(+\ TM+D3\) &
\(D2\) を transfer monotonicity へ弱める方向。\\
\(\exists\varphi\,T\nvdash(\Boxop\varphi\to\Dia\varphi)\) &
\(\Sigma_1 C\) &
Jeroslow 型。結論を \(\Con_T^L\) まで強めることは一般には不可。\\
\(\exists\varphi\,T\nvdash(\Boxop\varphi\to\Dia\varphi)\) &
\(M+D3\) &
Kurahashi 型。Rosser 述語との境界を示す。\\
\(T\nvdash\Con_T^S\) &
\(E+D3\) &
等値置換 \(E\) と introspection \(D3\) で strong consistency へ届く。\\
\(T\nvdash\Con_T^H\) &
\(M+CB+\Delta_0 C^U\) &
Hilbert--Bernays 型。全式量化を扱うため \(CB\) が現れる。\\
\(T\nvdash\Con_T^G\) &
\(D2^G+PC^G\) &
Montagna 型。global 条件を使う。\\
\(\Sigma_1 C^U\) と G2 系 & \(D1^U+D2^U\) &
Buchholz 型。uniform 条件から provable \(\Sigma_1\)-completeness へ。\\
\bottomrule
\end{longtable}

\section{\(\Gtwo\), \(\FGtwo\), 固定点の第一関係}

\begin{theorem}[Beklemishev--Shamkanov 型の抽象 G2]
\(\APS\) が A1--A4 を満たし，\(\Btx\)-固定点
\[
\exists p\in L,\quad p=\Btx p
\]
を持つとする。このとき \(\Gtwo\) と \(\FGtwo\) が成り立つ。
\end{theorem}

\begin{proof}
\(\Gtwo\). \(d=\Btx\Top\le\Bot\) とする。A2 と A1 の反単調性より
\[
\Top\le\Btx\Bot\le\Btx p=p.
\]
さらに \(\Top\le p\) から \(\Btx p\le\Btx\Top=d\) なので
\[
\Top\le p=\Btx p\le d\le\Bot.
\]

\(\FGtwo\). A4 より
\[
p=\Btx p\le\Boxop\Btx p=\Boxop p.
\]
また \(p\le\Btx p\) であるから，A3 により \(p\le\Btx\Top=d\)。
反単調性で
\[
\Btx d=\Btx\Btx\Top\le\Btx p=p\le d
\]
となる。
\end{proof}

\begin{lemma}[\(\FGtwo\Rightarrow\Gtwo\)]
A1 と A2 のもとで \(\FGtwo\) は \(\Gtwo\) を導く。
\end{lemma}

\begin{proof}
\(d=\Btx\Top\le\Bot\) とする。反単調性より \(\Btx\Bot\le\Btx d\)。
A2 と \(\FGtwo\) から
\[
\Top\le\Btx\Bot\le\Btx d\le d\le\Bot.
\]
\end{proof}

\begin{remark}[逆向きは反例で壊れる]
有限 preAPS zoo には \(\Gtwo\) だが \(\FGtwo\) でないモデル，\(\FGtwo\) だが
\(\Gtwo\) でないモデル，\(\Gtwo+\FGtwo\) だが syntactic fixed point を持たない
モデルが存在する。ただしどの APS axioms を保持するか，bottom discipline を課すか，
剰余付きにするかで witness は変わる。
\end{remark}

\section{Formalized G2 と local G2}

\(\Gtwo\) は外部 implication である一方，\(\FGtwo\) は「その implication を
\(\Btx\)-軌道内で内部化したもの」と読める。より細かい中間層として，
既存ノート \path{research/notes/formalized-g2-implicational-aps.md} では
local G2 が提案されている。

\begin{definition}[local G2]
\(d=\Btx\Top\) とする。各 \(a\in L\) に対し
\[
\mathrm{LG2}(a):\quad d\le a\Rightarrow \Btx a\le d
\]
と定義する。
\end{definition}

\begin{proposition}[局所化された鎖]
A1 のもとで
\[
\FGtwo
\Rightarrow
\forall a\,\mathrm{LG2}(a)
\Rightarrow
\mathrm{LG2}(\Bot)
\Rightarrow
\Gtwo
\]
が成り立つ。
\end{proposition}

\begin{proof}
\(\FGtwo\) と \(d\le a\) から，反単調性により
\[
\Btx a\le \Btx d\le d.
\]
よって \(\mathrm{LG2}(a)\)。\(\mathrm{LG2}(\Bot)\Rightarrow\Gtwo\) は，
\(d\le\Bot\) のとき A2 と \(\Btx\Bot\le d\) から
\[
\Top\le\Btx\Bot\le d\le\Bot
\]
となるため。
\end{proof}

\begin{remark}
\(\mathrm{LG2}(a)\) は order-only APS でも述べられるが，
\[
\Top\le (d\to a)\to(\Btx a\to d)
\]
のような単一命題として扱うには，implication や internal sequent object が必要である。
この違いが \(\Gtwo\) と \(\FGtwo\) の間にある「まだ名前のない中間層」である。
\end{remark}

\section{Löb と formalized Löb}

\begin{definition}[Löb rule と GL axiom]
順序的には Löb rule を
\[
\Boxop a\le a\quad\Longrightarrow\quad \Top\le a
\]
と読む。modal/provability logic では，formalized Löb は
\[
\Boxop(\Boxop\varphi\to\varphi)\to\Boxop\varphi
\]
すなわち GL axiom である。
\end{definition}

\begin{proposition}[Löb は固定点を consistency orbit に付着させる]
標準的な D1--D3/HBL 条件，あるいは GL のもとでは，
\[
g\leftrightarrow\neg\Boxop g
\]
で与えられる Gödel 型固定点は
\[
\neg\Boxop\Bot=\Con_T^L
\]
と同一視される。APS 的には，固定点が \(\Top\) から始まる consistency orbit の
第一段に付着する。
\end{proposition}

\begin{remark}[Rosser は脱付着として読む]
Rosser 型証明可能性述語は D1 や provable \(\Sigma_1\)-completeness を保ちながら，
formalized Löb/D3 の付着強制から逃げる。リポジトリ内の
\path{artifacts/reports/rosser-arithmetic-lift-check.json} は，
「monotonicity だけでは detached fixed point を排除できず，排除しているのは
formalized Löb/D3 である」ことを有限 GL frame と \(R_2\) 型順序で検証している。
\end{remark}

\subsection{軌道降下としての Löb}

有限順序では，\(\Top\)-軌道
\[
\Top,\Btx\Top,\Btx^2\Top,\ldots
\]
が最終的に安定化することを orbit descent と呼ぶ。検証済みレポート
\begin{quote}\small
\path{artifacts/reports/attachment-loeb-dividing-line-check.json}
\end{quote}
は，M3 型 carrier 上で次を確認している。

\begin{center}
\begin{tabularx}{\linewidth}{>{\raggedright\arraybackslash}p{0.27\linewidth}
>{\raggedright\arraybackslash}X}
\toprule
有限検証結果 & 読み方 \\
\midrule
descending orbit \(\Rightarrow\) attached fixed point &
軌道が止まるなら，その停止値は orbit-attached な固定点である。これは de Jongh--Sambin/Löb 側の影。\\
detached-only fixed point \(\Rightarrow\) non-descent \(+\ \neg\FGtwo\) &
Rosser 型の detached fixed point は，\(\Top\)-軌道の非降下・循環として現れる。\\
\(\neg\FGtwo\) だけでは detachment を強制しない &
\(\FGtwo\) 失敗後に attached fixed point へ到達するモデルもある。真の境界は orbit descent。\\
\bottomrule
\end{tabularx}
\end{center}

\section{不動点 zoo}

\subsection{四つの代表型}

\begin{longtable}{>{\raggedright\arraybackslash}p{0.17\linewidth}
>{\raggedright\arraybackslash}p{0.27\linewidth}
>{\raggedright\arraybackslash}p{0.43\linewidth}}
\toprule
名前 & 典型方程式 & G2/APS 上の読み方 \\
\midrule
\endfirsthead
\toprule
名前 & 典型方程式 & G2/APS 上の読み方 \\
\midrule
\endhead
Gödel 文 &
\(g\leftrightarrow\neg\Boxop g\) &
\(\Btx=\neg\Boxop\) と定義できる古典的環境では \(\Btx\)-固定点。
GL/Löb の下では consistency orbit に付着する。\\
Henkin 文 &
\(h\leftrightarrow\Boxop h\) &
「自分は証明可能」と言う文。Löb により，Henkin 文であることが証明されれば
\(T\vdash h\) となる。\\
Santa Claus 文 &
\(s_\psi\leftrightarrow(\Boxop s_\psi\to\psi)\) &
パラメータ \(\psi\) に対する固定点族。\(\forall\exists SC\) は強い固定点存在原理だが，
それ単独では Löb を含意しない。\\
Jeroslow 文 &
\(j\leftrightarrow\Boxop\neg j\), または \(j=\Btx j\) &
primitive \(\Btx\) を許すと，Gödel 型より弱い構造で存在しうる。
BS16/APS では fixed point が \(\Gtwo,\FGtwo\) を導くが，contraction-free な剰余付き読解では
formalized G2 と分離される。\\
\bottomrule
\end{longtable}

\begin{remark}[primitive refutability と negated provability を混同しない]
既存ノート
\begin{quote}\small
\path{research/notes/mnd4-preaps-fixedpoint-obstruction.md}
\end{quote}
の重要点は
\[
p=\Btx p
\quad\text{と}\quad
p=\neg\Boxop p
\]
が別物だということである。古典否定，爆発，D/4 型 introspection，十分な contraction を
入れると後者は崩壊に近い力を持つ。一方，primitive \(\Btx\)-固定点は，
preAPS や substructural setting では非自明に存在できる。
\end{remark}

\subsection{Santa Claus は Löb を含意しない}

既存ノート \path{research/notes/shibuya-seminar-2026-05-08.md} の連続区間例は，
\(\exists SC\not\Rightarrow\) Löb を示す。

\[
L=[0,1],\qquad
x\to y=\min(1,1-x+y)
\]
を Łukasiewicz implication とし，\(\Boxop=\id\) とする。このとき
\[
\varphi=\Boxop\varphi\to\psi
\]
は \(\varphi=\varphi\to\psi\) であり，右辺は連続写像なので中間値の定理から固定点が存在する。
したがって \(\exists SC\) は成り立つ。一方 Löb は
\[
\Boxop a\le a\Rightarrow \Top\le a
\]
だが，\(\Boxop=\id\) なので前提は全ての \(a\) で真であり，結論は
\(a<1\) で偽となる。したがって Löb は失敗する。

\section{Cut elimination zoo と剰余付き APS}

\subsection{A3 は隠れた contraction}

BS16 型の contraction-free provability logic を剰余付き preAPS として読むと，
\[
\Btx A:=\Boxop(A\to\Bot)
\]
であり，
\[
\FGtwo:\quad \Btx\Btx\Top\le\Btx\Top
\]
は
\[
\Boxop(\Boxop\Bot\to\Bot)\le \Boxop\Bot
\]
になる。これは formalized G2 の形そのものである。

一方，APS A3
\[
x\le\Boxop y,\quad x\le\Btx y\Rightarrow x\le\Btx\Top
\]
は，資源敏感に読むと，二つの仮定を同じ \(x\) から同時に使っている。
剰余付き世界では自然には
\[
x\otimes x\le\Boxop y\otimes\Btx y
\]
となり，そこから \(x\) に戻るには diagonal contraction 的な原理が必要である。
したがって
\[
\text{A3}
=
\text{provability/refutability collision}
+\text{hidden contraction}
\]
と読むのが自然である。

\subsection{cut elimination の読解}

cut-free proof values を fiber
\[
\mathcal E^{\cf}_{\Gamma,\Delta}
\]
として持つとき，cut elimination は
\[
\pi_1\in\mathcal E^{\cf}_{\Gamma;\Delta,A},\quad
\pi_2\in\mathcal E^{\cf}_{A,\Sigma;\Pi}
\Rightarrow
\pi_2\circ_{\cut}\pi_1\in
\mathcal E^{\cf}_{\Gamma,\Sigma;\Pi,\Delta}
\]
へ正規化できるという安定性である。代数的には，これは MacNeille/canonical completion
に対する閉性，または residuated completion に対する安定性として現れる。

\begin{center}
\begin{tabularx}{\linewidth}{>{\raggedright\arraybackslash}p{0.25\linewidth}
>{\raggedright\arraybackslash}p{0.31\linewidth}
>{\raggedright\arraybackslash}X}
\toprule
proof-theoretic 側 & algebraic 側 & G2 zoo で見る量 \\
\midrule
cut elimination & cut-free order の合成閉性 & MacNeille/canonical completion の安定性 \\
contraction なし & non-idempotent \(\otimes\) & fixed point はあっても \(\FGtwo\) が失敗しうる \\
restricted contraction & 局所的な \(a\otimes a\le a\) & front ideal の局所 contraction \\
full contraction & idempotent/Heyting 的構造 & A3 と formalized G2 が復活しやすい \\
weakening & \(a\le b\Rightarrow a\otimes c\le b\) & bottom discipline，discarding rule \\
exchange & \(a\otimes b=b\otimes a\) & 現在の \(B_N\) residuated witnesses では多くが満たす \\
\bottomrule
\end{tabularx}
\end{center}

\section{抽象的導出可能性条件}

本稿で使う主要条件をまとめる。

\begin{longtable}{>{\raggedright\arraybackslash}p{0.13\linewidth}
>{\raggedright\arraybackslash}p{0.47\linewidth}
>{\raggedright\arraybackslash}p{0.27\linewidth}}
\toprule
記号 & 形式 & 役割 \\
\midrule
\endfirsthead
\toprule
記号 & 形式 & 役割 \\
\midrule
\endhead
\(C\) &
\(\Boxop(\varphi\to\psi)\to(\Boxop\varphi\to\Boxop\psi)\) &
K axiom，internal modus ponens。\\
\(M\) &
\(\varphi\to\psi\) ならば \(\Boxop\varphi\to\Boxop\psi\) &
local monotonicity。\\
\(E\) &
\(\varphi\leftrightarrow\psi\) ならば \(\Boxop\varphi\leftrightarrow\Boxop\psi\) &
extensionality，等値置換。\\
\(D3\) &
\(\Boxop\varphi\to\Boxop\Boxop\varphi\) &
introspection，modal \(4\)。\\
\(\Sigma_1 C^U\) &
\(\theta(x)\to\Boxop\theta(\dot x)\) for \(\Sigma_1\ \theta\) &
Jeroslow/Kurahashi 型の入口。\\
\(\Delta_0 C^U\) &
\(\theta(x)\to\Boxop\theta(\dot x)\) for \(\Delta_0\ \theta\) &
Hilbert--Bernays 型で使う。\\
\(CB\) &
\(\Boxop\forall x\,\varphi(x)\to\forall x\,\Boxop\varphi(\dot x)\) &
全称文の box collection。\\
\(PC^G\) &
\(\forall x(Pr_{\emptyset}(x)\to Pr_T(x))\) &
global predicate calculus completeness。\\
\(TM\) &
\(\varphi\to(\psi\to\gamma)\) ならば
\(\Boxop\varphi\to(\Boxop\psi\to\Boxop\gamma)\) &
\(D2\) の transfer 型弱化。\\
\(\Ros\) &
\(\neg\varphi\Rightarrow\neg\Boxop\varphi\) &
Rosser rule，detached fixed point 側。\\
\(\mathrm{Loeb}^U\) &
\(\forall x(\Boxop\varphi(x)\to\varphi(x))\Rightarrow\forall x\,\varphi(x)\) &
uniform Löb rule。\\
\bottomrule
\end{longtable}

\begin{proposition}[条件間の基本関係]
標準的な読みでは
\[
D2\quad\Longleftrightarrow\quad M+C
\]
であり，また
\[
\mathrm{Global}\Rightarrow\mathrm{Uniform}\Rightarrow\mathrm{Local}
\]
という強さの方向がある。ただし \(M,TM,wM\) の細部は，Rosser 述語や non-normal
modal logic の例で分離されうる。
\end{proposition}

\section{反例モデル表}

\subsection{有限モデル zoo}

以下は \Repo{} 内の machine-checked reports から，今回の関係整理に必要なものだけを
抜き出した表である。

\begin{longtable}{>{\raggedright\arraybackslash}p{0.19\linewidth}
>{\raggedright\arraybackslash}p{0.18\linewidth}
>{\raggedright\arraybackslash}p{0.25\linewidth}
>{\raggedright\arraybackslash}p{0.25\linewidth}}
\toprule
モデル & 構造クラス & 成立/失敗 & 何を分離するか \\
\midrule
\endfirsthead
\toprule
モデル & 構造クラス & 成立/失敗 & 何を分離するか \\
\midrule
\endhead
\texttt{M-010} &
3-element preAPS &
\(\FGtwo\) true, \(\Gtwo\) false, no FP-synt &
\(\FGtwo\not\Rightarrow\Gtwo\) の小 witness。ただし bottom discipline 修復後はプロファイルが変わる。\\
\texttt{M-100} &
3-element preAPS &
\(\Gtwo\) true, \(\FGtwo\) false, no FP-synt &
\(\Gtwo\not\Rightarrow\FGtwo\)。\\
\texttt{M-101} &
3-element preAPS &
\(\Gtwo\) true, \(\FGtwo\) false, FP-synt true &
\(\Gtwo+\) fixed point でも \(\FGtwo\) へ届かないことを，A3/A4 パッケージ外で示す。\\
\texttt{M-110} &
3-element preAPS &
\(\Gtwo\) true, \(\FGtwo\) true, FP-synt false &
\(\Gtwo+\FGtwo\not\Rightarrow\) syntactic fixed point。\\
\texttt{bottom-nfg2-depth-5-front-shifted-non-u-absorbing} &
bottom-disciplined full-residuated preAPS &
\(\Gtwo\) true, \(\FGtwo\) false, first true \(\nFGtwo\) depth \(6\), FP-synt \(s\) &
\(\nFGtwo\) 階層を任意深く遅らせられる。front-shifted tensor は \(U\)-absorbing でない。\\
\texttt{bottom-G2FG2-noFP-residuated} &
bottom-disciplined full-residuated preAPS &
\(\Gtwo\) true, \(\FGtwo\) true, FP-synt false &
剰余付きでも \(\Gtwo+\FGtwo\) が fixed point を強制しない。\\
\texttt{R\_2} gadget &
M3 diamond order &
detached FP \(p\), orbit \(o_0\leftrightarrow o_1\), \(\FGtwo\) false &
Rosser 文の代数的影。固定点が consistency orbit から完全に外れる。\\
3-chain attached model &
chain/order-attached model &
orbit descends to fixed point &
Löb/de Jongh--Sambin 側。固定点が consistency orbit に付着する。\\
\([0,1]\) Łukasiewicz SC model &
implication-extended continuous ABS &
\(\exists SC\) true, Löb false &
Santa Claus fixed points do not imply Löb。\\
\bottomrule
\end{longtable}

\subsection{無限モデルまたは無限 family}

\begin{enumerate}
\item \textbf{\(B_N\) family.}
リポジトリの \texttt{bottom-nfg2-depth-\(N\)} 系列は，任意の有限深さまで
\(\nFGtwo\) の初回成立を遅らせる。したがって単一の有限 bound では
\(\FGtwo\) と全階層 \(\nFGtwo(k)\) の差を潰せない。

\item \textbf{period-\(2k\) Rosser gadgets.}
M3 diamond の \(R_2\) を \(2k\)-周期の antichain orbit に拡張すると，
detached fixed point と非降下 orbit を保ったまま周期を任意に大きくできる。
これは Rosser 側の「偶数周期 gap」の有限 family である。

\item \textbf{infinite antichain Rosser template.}
ABS/preAPS レベルなら，
\[
L=\{\Bot,\Top,p\}\cup\mathbb Z
\]
で \(\Bot\) を最小，\(\Top\) を最大，\(\mathbb Z\cup\{p\}\) を antichain とし，
\[
\Btx\Bot=\Top,\quad \Btx\Top=\Bot,\quad
\Btx n=n+1,\quad \Btx p=p
\]
と置くことで，detached fixed point \(p\) と無限に降下しない consistency orbit を
同時に持つ無限 \(\ABS\) が得られる。これは A3 を満たす \(\APS\) witness ではなく，
Rosser geometry だけを取り出すための反例テンプレートである。
\end{enumerate}

\section{完成と固定点の反映問題}

MacNeille completion などの completion は固定点を作ることがある。しかし本プロジェクトの
基本 slogan は
\[
\text{completion fixed point}
\neq
\text{syntactic fixed point}
\]
である。

\begin{definition}[MacNeille completion]
\[
X^u=\{a:\forall x\in X,\ x\le a\},\qquad
X^l=\{a:\forall x\in X,\ a\le x\}.
\]
MacNeille completion \(\widehat L\) は \(C=(C^u)^l\) を満たす closed lower cut 全体であり，
principal embedding は
\[
i(a)=(\{a\}^u)^l
\]
である。
\end{definition}

反単調な \(\Btx:L\to L\) は，まず \(L\to L^{op}\) の単調写像として extension すべきである。
現在の checker convention は
\[
\widehat{\Btx}(C)=\bigl((\Btx[C])^{l_L}\bigr)^{u_L}
\]
であり，これは \(L^{op}\)-MacNeille closure として読む。

\begin{center}
\begin{tabularx}{\linewidth}{>{\raggedright\arraybackslash}p{0.25\linewidth}
>{\raggedright\arraybackslash}X}
\toprule
completion FP の型 & 意味 \\
\midrule
reflected principal & \(q=i(p)\) かつ \(p=\Btx p\)。構文的固定点へ戻る。\\
principal-unreflected & \(q=i(a)\) だが \(a\ne\Btx a\)。主成分だが固定点ではない。\\
non-principal & cut/ideal としてのみ存在し，式としては戻らない。\\
none & completion 後にも固定点なし。\\
\bottomrule
\end{tabularx}
\end{center}

有限検証では，\texttt{bottom-G2FG2-noFP-residuated} などが
principal-unreflected completion fixed point を持つ。これは，\(\Gtwo+\FGtwo\) と
completion fixed point があっても，syntactic fixed point が戻るとは限らないことを示す。

\section{全体関係図}

現在安全に使える関係を，条件付きの矢印としてまとめる。
\[
\exists p(p=\Btx p)+\mathrm{A1\text{--}A4}
\Longrightarrow
\FGtwo
\Longrightarrow
\forall a\,\mathrm{LG2}(a)
\Longrightarrow
\mathrm{LG2}(\Bot)
\Longrightarrow
\Gtwo.
\]
\[
\mathrm{GL/Loeb}
\Longrightarrow
\text{Gödel fixed point is orbit-attached}
\Longrightarrow
\text{Rosser }R_2\text{ geometry is impossible}.
\]
\[
\begin{aligned}
&\text{Rosser/detached fixed point}\\
&\qquad\Longrightarrow
\text{failure of formalized Loeb/D3 package}\\
&\qquad\text{but not necessarily failure of monotonicity}.
\end{aligned}
\]
\[
\begin{aligned}
&\text{cut elimination in a residuated/fibered calculus}\\
&\qquad\Longleftrightarrow
\text{closure of cut-free order under cut composition}\\
&\qquad\Longrightarrow
\text{completion stability questions for }\Gtwo,\FGtwo,\Fix.
\end{aligned}
\]
\[
\exists SC\not\Rightarrow\mathrm{Loeb},
\qquad
\Gtwo+\FGtwo\not\Rightarrow\exists p(p=\Btx p).
\]
The last separation is witnessed in the finite preAPS/RAPS zoo.

\section{未解決問題}

\begin{enumerate}
\item \(\Con_T^L\), \(\Ros\), \(\Con_T^S\), \(\Con_T^H\) の間を，導出可能性条件
\(C,M,E,D3,CB,\Sigma_1C\) の最小パッケージで完全分類する。
\item \(\Con_T^{IG}\) と \(\Con_T^{EG}\) を，APS 内部の可算 join，式コード object，
または indexed/fibred APS として厳密に分離する。
\item \(\mathrm{LG2}(a)\) と \(\FGtwo[q]\) を finite checker に入れ，
\[
\FGtwo\Rightarrow \forall a\,\mathrm{LG2}(a)\Rightarrow\mathrm{LG2}(\Bot)\Rightarrow\Gtwo
\]
の各矢印の strictness を検証する。
\item BS16 cut elimination を，fibered residuated APS の MacNeille/canonical completion
安定性として定式化する。
\item completion fixed point が syntactic fixed point へ反映されるための compactness，
definability，または diagonalization 条件を抽出する。
\item 無限 \(B_\omega\) 型または \(\omega\)-orbit 型 RAPS で，全ての有限 \(\nFGtwo(k)\)
を失敗させる剰余付き witness を構成できるか調べる。
\end{enumerate}

\section*{参照したローカル成果物}

\begin{itemize}
\item \path{research/definitions.md}
\item \path{research/notes/consistency-notions-hierarchy.md}
\item \path{research/notes/derivability-conditions.md}
\item \path{research/notes/g2-zoo-arithmetic.md}
\item \path{research/notes/formalized-g2-implicational-aps.md}
\item \path{research/notes/g2-aps-zoo-classification.md}
\item \path{research/notes/bs16-fiber-residuated-aps.md}
\item \path{research/notes/mnd4-preaps-fixedpoint-obstruction.md}
\item \path{research/notes/completion-and-fixed-points.md}
\item \path{research/notes/shibuya-seminar-2026-05-08.md}
\item \path{artifacts/reports/detached-rosser-fixedpoint-check.json}
\item \path{artifacts/reports/attachment-loeb-dividing-line-check.json}
\item \path{artifacts/reports/rosser-arithmetic-lift-check.json}
\item \path{artifacts/reports/g2-zoo-bottom-nfg2-depth-5-front-shifted-non-u-absorbing.json}
\item \path{artifacts/reports/g2-zoo-bottom-G2FG2-noFP-residuated.json}
\end{itemize}

\section*{外部文献}

\begin{itemize}
\item L. Beklemishev and D. Shamkanov,
``Some abstract versions of Gödel's second incompleteness theorem based on non-classical logics,''
\emph{Logic, Methods and Philosophy of Science}, 12(2), 2016.
\href{https://arxiv.org/abs/1602.05728}{arXiv:1602.05728}.
\item T. Kurahashi,
``A note on derivability conditions,''
\href{https://arxiv.org/abs/1902.00895}{arXiv:1902.00895}.
\item T. Kurahashi,
``Rosser provability and the second incompleteness theorem,''
\href{https://arxiv.org/abs/1902.06863}{arXiv:1902.06863}.
\item T. Kurahashi,
``Refinements of provability and consistency principles for the second incompleteness theorem,''
\href{https://arxiv.org/abs/2507.00955}{arXiv:2507.00955}.
\item R. G. Jeroslow,
``Redundancies in the Hilbert-Bernays derivability conditions for Gödel's second incompleteness theorem,''
\emph{Journal of Symbolic Logic}, 1973.
\item D. Hilbert and P. Bernays, \emph{Grundlagen der Mathematik II}, 1939.
\item M. H. Löb,
``Solution of a problem of Leon Henkin,''
\emph{Journal of Symbolic Logic}, 1955.
\item C. Smoryński, \emph{Self-Reference and Modal Logic}, Springer, 1985.
\item G. Boolos, \emph{The Logic of Provability}, Cambridge University Press, 1993.
\item D. Guaspari and R. Solovay,
``Rosser sentences,''
\emph{Annals of Mathematical Logic}, 16, 1979.
\item N. Galatos, P. Jipsen, T. Kowalski, and H. Ono,
\emph{Residuated Lattices: An Algebraic Glimpse at Substructural Logics},
Elsevier, 2007.
\item A. Ciabattoni, N. Galatos, and K. Terui,
work on algebraic proof theory, analytic rules, cut elimination, and completions for substructural logics.
\item S. Iwata and T. Kurahashi,
``Fixed-point properties for predicate modal logics,''
\href{https://arxiv.org/abs/1907.00306}{arXiv:1907.00306}.
\end{itemize}

\end{document}
